\section{Introduction}

Synthetic financial data generation serves critical functions in quantitative finance: stress testing portfolios against scenarios that have never occurred, estimating risk metrics with larger sample sizes than history provides, and preserving privacy when sharing proprietary trading data.
The challenge is acute for multi-asset portfolios, where generating realistic cross-asset dependence structures---correlations, tail co-movement, and volatility spillovers---is substantially harder than matching marginal distributions of individual assets.

Traditional approaches fall into two camps.
Parametric models such as DCC-GARCH \citep{engle2002dynamic} impose strong distributional assumptions that often fail to capture the nonlinear, regime-dependent dynamics of real markets.
Non-parametric methods like factor bootstrap \citep{huh2026marketgans} preserve the empirical structure of historical data but cannot generate counterfactual scenarios conditioned on hypothetical macro states.
Recent work on financial GANs, including MarketGAN \citep{huh2026marketgans} and related architectures \citep{vuletic2023fingen, cont2025tailgan}, attempts to bridge this gap by learning the full return-generating process adversarially.
These models typically learn corrections to factor loadings ($\beta$), idiosyncratic volatilities ($\sigma$), and residual dynamics ($\varepsilon$) simultaneously.

We find, however, that this full-learning approach systematically hurts OOS generalization.
Ablations show two consistent patterns:
(1)~learned residual networks inject spurious cross-class correlation that does not persist out of sample, and
(2)~adversarial corrections to factor loadings distort the cross-asset covariance structure that rolling OLS already captures well.
A third pattern emerges from a focused failure-mode study: purely data-driven (PCA) factor decompositions leave \emph{same-sector} pair co-movement in the residual space, and any fixed-Cholesky residual prior then becomes a bottleneck for those pairs specifically. For instance, with only global + within-class PCA factors, the Visa--Mastercard training-period residual correlation is $+0.57$ (out-of-sample $+0.73$); a shrinkage-toward-identity residual prior pushes this to $+0.17$, catastrophically underestimating the payments co-movement that a risk analyst cares most about.

These observations motivate two coupled design choices: (i) a \emph{sector-augmented, Gram--Schmidt-orthogonalised} factor set that puts the same-sector common component into $\betahat$ rather than the residual; and (ii) a \emph{factor-structured generator} that restricts the adversarial component to quantities where classical estimators are demonstrably weak.

\paragraph{Our contributions.}
\begin{itemize}
    \item \textbf{Sector-augmented factor decomposition (representation contribution).}
    We extend the hierarchical PCA factor model of our prior MarketGAN work with 5 custom Fama--French-style stock industry factors (splitting \textsc{Finance} out of FF5's ``Other'' so Visa/Mastercard/JPM/BRK.B share a bucket) and 4 S\&P GSCI commodity sub-indices (\textsc{Energy}, \textsc{Precious}, \textsc{Industrial}, \textsc{Agriculture}), then apply sequential Gram--Schmidt orthogonalisation and unit-variance standardisation. The resulting $K=16$ factor panel has training off-diagonal correlation $0$ and $\beta$-magnitudes a factor of $40$ smaller than the raw (non-standardised) equal-weight sector means, making rolling OLS numerically stable at $N/K \approx 3.8$. Ablations (\S\ref{sec:ablation}) show this alone reduces OOS correlation Frobenius by $11\%$ relative to the 11-factor PCA-only baseline.

    \item \textbf{Factor-structured generator (architectural contribution).}
    We restrict generator capacity to macro-conditioned corrections on intercepts ($\alpha$) and idiosyncratic volatilities ($\sigma$) via a TCN backbone with FiLM conditioning. Factor loadings ($\beta$) are taken directly from rolling OLS. Residuals use a \emph{fixed} training-sample correlation Cholesky $L_\rho$ (not a shrinkage prior, not a learned network); Appendix~\ref{app:L_rho_ablation} reports three failed attempts to make $L_\rho$ adversarially trainable.

    \item \textbf{Regime-gated correction schedule (mechanism contribution).}
    A VIX-aligned encoder gates the magnitudes of the $(\alpha, \sigma)$-correction heads: a tight bound ($\eta_\alpha = 0.01$, $\eta_\sigma = 0.05$) in calm regimes, relaxed to $(\eta_\alpha = 0.10$, $\eta_\sigma = 0.20)$ under stress. The encoder is trained with an auxiliary BCE loss aligned to a hard VIX-z threshold, which we verify with a hard-threshold ablation.

    \item \textbf{Main 60-asset OOS.}
    Under strictly lagged factor conditioning ($f_{t-1}$), SF-MarketGAN-R achieves OOS correlation Frobenius $2.53$ vs.\ factor-bootstrap $3.36$ ($-24.6\%$) and shrinkage-prior baseline $3.12$ ($-18.8\%$), winning $4$ of $6$ evaluation metrics. The three stylized-fact metrics (ACF, vol-clustering, leverage) each improve by $10$--$13\%$ over the same-factor FB+shrunk baseline---a structural advantage that a better residual Cholesky alone cannot produce.

    \item \textbf{Same-sector pair-correlation baseline comparison.}
    We benchmark against $7$ classical and neural alternatives on six same-sector pairs covering payments, precious metals, Treasuries, energy, mega-tech, and cross-subgroup. Factor-bootstrap with an appropriately estimated residual Cholesky matches SF-MarketGAN-R on all six pairs within measurement noise---so the GAN's incremental value is \emph{not} on pair fidelity but on structural stylized-fact fit that no classical method approaches.
\end{itemize}

The rest of this paper is organised as follows. Section~\ref{sec:related} reviews related work. Section~\ref{sec:method} presents the model architecture and its theoretical motivation. Section~\ref{sec:experiments} reports experimental results including the pair-correlation baseline study, stylized-fact comparison, and ablations. Section~\ref{sec:conclusion} discusses limitations and future directions.
